\begin{textsample}{POS Dim 4 – 2017\_c5 – Score 1.00 – 2017\_c5/enade\_2017\_eng\_qui\_c6961b35ac.txt}  \label{ex:f4_pos_229}
CAPÍTULO 5 DISTRIBUIÇÃO DOS CONCEITOS 5.1 PANORAMA NACIONAL DA DISTRIBUIÇÃO DOS CONCEITOS O Conceito Enade é um indicador de qualidade que avalia os cursos por intermédio dos desempenhos dos estudantes a partir dos resultados do Enade.
Os valores possíveis vão de 1 ( pior situação ) a 5 ( melhor situação ).
Os cursos com apenas um ou sem nenhum concluinte participante não obtêm o Conceito Enade, ficando Sem Conceito ( SC ).
A Tabela 5.1 apresenta a quantidade e distribuição de cursos de Engenharia Química participantes do Enade/2017, por faixa de conceito e Grande Região.
Enfatiza -se, mais uma vez, que a diferença entre os cursos tabulados neste capítulo e no capítulo 2 corresponde aos cursos Sem Conceito ( SC ), ou seja àqueles, em princípio, sem alunos concluintes que houvessem participado da prova ou àqueles que tivessem tão somente um aluno concluinte18.
Observando -se os dados da Tabela 5.1, nota -se que, dos 135 cursos participantes, 50 ( 37,0 \% ) classificaram -se com conceito 3, o valor modal.
Esse foi também o conceito modal na região Norte ( 60,0 \%, correspondendo a 3 cursos ) e na região Sul ( 44,7 \%, correspondendo a 17 cursos ).
Na região Centro-Oeste, a moda ficou dividida entre os conceitos 3 e 4 ( um curso, correspondendo a 50,0 \%, cada um ).
Na região Sudeste ( 35,2 \%, correspondendo a 25 cursos ), o conceito modal foi 2.
Já na região Nordeste ( 36,8 \%, correspondendo a 7 cursos ), o conceito modal foi 4.
O conceito 2 foi o segundo mais frequente em nível nacional ( 28,1 \%, correspondendo a 38 cursos ), e o conceito 4, o terceiro ( 25,2 \%, correspondendo a 34 cursos ).
Houve, ainda, nove cursos ( 6,7 \% ) que receberam conceito 5, e dois cursos ( 1,5 \% ), que receberam conceito 1.
Dos 135 cursos de Engenharia Química, dois ( 1,5 \% ) ficaram Sem Conceito ( SC ). 18 Esses cursos com somente um concluinte não foram considerados no capítulo 2.
A região Norte participou com cinco cursos ou 3,7 \% do total nacional.
Desses, três cursos receberam o conceito 3, o que equivale a 60,0 \% do total regional, além de ser o conceito modal na região, como já comentado.
Os conceitos 2 e 4 foram atribuídos a um curso, cada um ( 20,0 \% para cada ).
Nenhum curso ficou Sem Conceito ( SC ), e nenhum curso recebeu os conceitos 1 ou 5 nessa região.
A região Nordeste participou com 19 cursos ou 14,1 \% do total nacional.
Nessa região, sete cursos ( 36,8 \% em termos regionais ) obtiveram conceito 4, o conceito modal para a região, como já comentado.
Os conceitos 2 e 3 foram atribuídos a cinco cursos cada ( 26,3 \%, cada ).
Ainda nessa região, dois cursos ( 10,5 \% ) foram avaliados com conceito 5, e nenhum curso ficou com conceito 1 ou Sem Conceito ( SC ).
Dos 71 cursos participantes da região Sudeste ( 52,6 \% do total nacional ), o conceito 2 foi o conceito modal, como já anteriormente assinalado, tendo sido obtido por 25 cursos ( 35,2 \% ).
O conceito 1 foi atribuído a dois cursos ( 2,8 \% ), e o conceito 3, a 24 cursos ( 33,8 \% ).
Receberam os conceitos 4 e 5, respectivamente, 13 cursos ( 18,3 \% ) e seis cursos ( 8,5 \% ).
Nessa região, um curso ( 1,4 \% ) ficou Sem Conceito ( SC ).
Os 38 cursos da região Sul corresponderam a 28,1 \% do total nacional.
Como já indicado, o conceito 3 foi o conceito modal, atribuído a 44,7 \% dos cursos da região, ou seja, atribuídos a 17 cursos.
O conceito 2 foi atribuído a sete cursos ( 18,4 \% ), o conceito 4 a 12 cursos ( 31,6 \% ), e o conceito 5, a um curso ( 2,6 \% ).
Um curso ficou Sem Conceito ( SC ), e nenhum curso recebeu conceito 1 na região Sul.
Já dos dois cursos participantes na região Centro-Oeste ( 1,5 \% do total nacional ), um deles obteve conceito 3, e o outro, conceito 4, o que corresponde a 50,0 \% para cada um.
Nessa região, nenhum curso ficou Sem Conceito ( SC ), assim como nenhum curso recebeu os demais conceitos.
O Gráfico 5.1 apresenta a distribuição cumulativa do Conceito Enade segundo a Grande Região.
Em um gráfico de distribuição cumulativa, em cada valor de abscissa ( eixo x ), o valor apresentado no eixo das ordenadas ( eixo y ) é o valor acumulado da variável.
Por exemplo, a região Norte ( linha azul ) apresenta 80 \% dos cursos com conceito 3 ou abaixo desse valor.
A região Nordeste ( linha verde ), por outro lado, apresenta um pouco mais de 52 \% dos cursos na mesma situação : com conceito 3 ou abaixo.
Na região Norte ( linha azul ) todos os cursos têm conceitos 4 ou abaixo disso, já que a poligonal cumulativa alcança 100 \% para o conceito 4.
A região Centro-Oeste ( linha preta ) também apresenta a mesma característica, diferenciando -se apenas pelo fato de não ter cursos com os conceitos 1 ou 2, já que a poligonal cumulativa parte de 0 \% no conceito 2.
Em linhas gerais, em um gráfico de distribuição cumulativa, poligonais mais à esquerda representam Grandes Regiões com pior distribuição, e poligonais mais à direita, Grandes Regiões com melhores conceitos.
A região Norte ( linha azul ) apresenta os cursos com conceitos mais baixos, ao passo que as regiões Nordeste ( linha verde ) e Centro-Oeste ( linha preta ) apresentam as distribuições com valores maiores, porém a região Centro-Oeste apresenta apenas dois cursos, como já mencionado.
Cursos em IES situadas nas regiões Sudeste ( linha vermelha ) e Sul ( linha roxa ) apresentam uma situação intermediária. 5.2 CONCEITOS POR CATEGORIA ADMINISTRATIVA, POR MODALIDADE DE ENSINO E GRANDE REGIÃO A Tabela 5.2 apresenta a distribuição dos cursos de Engenharia Química participantes do Enade/2017, por Categoria Administrativa e Modalidade de ensino, de acordo com os conceitos por eles alcançados, segundo as Grandes Regiões brasileiras.
Dos 135 cursos participantes, 61 ( 45,2 \% ) eram ministrados em Instituições Públicas, e 74 ( 54,8 \% ), em Instituições Privadas.
Quanto à modalidade de ensino, todos os 135 cursos eram presenciais.
De acordo com as informações da Tabela 5.2, em termos nacionais, entre Instituições Públicas e Privadas, observa -se que, dos nove cursos avaliados com conceito 5, todos eram oferecidos em IES Públicas.
Dos 61 cursos participantes de IES Públicas, o conceito 4 foi o valor modal, atribuído a 27 cursos ( 44,3 \% ).
Dentre os demais cursos de Instituições Públicas participantes, três cursos ( 4,9 \% ) foram avaliados com conceito 2, e 20 cursos ( 32,8 \% ), com conceito 3.
Como já comentado, nove cursos ( 14,8 \% ) receberam conceito 5.
Nessa categoria, dois cursos ( 3,3 \% ) ficaram Sem Conceito ( SC ), e nenhum ficou com conceito 1.
Na rede privada, o conceito modal foi 2, atribuído a 35 cursos, o correspondente a 47,3 \% dos 74 cursos da categoria.
Dentre os demais cursos participantes, dois ( 2,7 \% ) receberam conceito 1, e 30 ( 40,5 \% ), conceito 3.
O conceito 4 foi atribuído a sete cursos ( 9,5 \% ).
Nessa Categoria Administrativa, nenhum curso ficou Sem Conceito ( SC ) ou recebeu conceito 5.
Como todos os cursos de Engenharia Química são presenciais, a configuração conceitual é a mesma descrita para o cenário nacional.
Na análise por região, observa -se que, na região Norte, as Instituições Privadas participaram com apenas um curso ( 20,0 \% do total regional ), que obteve conceito 2.
As Instituições Públicas participaram com quatro cursos na região Norte ( 80,0 \% do total regional ), dos quais três ( 75,0 \% ) obtiveram conceito 3, o valor modal.
O conceito 4 foi atribuído a um curso ( 25,0 \% ).
Nenhum curso ficou Sem Conceito ( SC ), assim como nenhum recebeu os demais conceitos nessa combinação de categoria e região.
Na região Nordeste, a rede privada concentrou seis dos 19 cursos participantes, o equivalente a 31,6 \% do total da região.
Dentre os cursos oferecidos por IES Privadas no Nordeste, cinco cursos ( 83,3 \% ) foram avaliados com conceito 2, o conceito modal.
Um curso obteve conceito 3 ( 16,7 \% ).
Nenhum curso ficou Sem Conceito ( SC ) ou com os demais conceitos nessa região.
As Instituições Públicas da região Nordeste participaram com 13 cursos ( 68,4 \% do total da região ).
Desses, sete obtiveram conceito 4 ( 53,8 \% ), o valor modal, seguidos por quatro cursos, que receberam conceito 3 ( 30,8 \% ), e por dois cursos, que receberam conceito 5 ( 15,4 \% ).
Na região Sudeste, a proporção de cursos da rede privada, 66,2 \%, foi mais elevada do que a registrada nas demais regiões brasileiras, correspondendo a 47 dos 71 cursos participantes.
Nessa categoria e região, o conceito modal foi 2, atribuído a 24 cursos ( 51,1 \% ).
Nenhum curso ficou Sem Conceito ( SC ) ou recebeu conceito 5.
Os demais foram avaliados com conceito 1 ( dois cursos, 4,3 \% ), conceito 3 ( 17 cursos, 36,2 \% ) e conceito 4 ( quatro cursos, 8,5 \% ).
Dentre os 24 cursos oferecidos por Instituições Públicas na região Sudeste, o conceito modal foi 4, atribuído a nove cursos ( 37,5 \% ).
Os demais cursos receberam os conceitos 2 ( um curso, 4,2 \% ), 3 ( sete cursos, 29,2 \% ) e 5 ( seis cursos, 25,0 \% ).
Na região Sudeste, nessa categoria, um curso ( 4,2 \% ) ficou Sem Conceito ( SC ), e nenhum curso recebeu conceito 1.
As Instituições Privadas concentraram 20 dos 38 cursos participantes da região Sul, 52,6 \% do total regional.
Desses, 12 ( 60,0 \% ) obtiveram conceito 3, o conceito modal.
Nessa combinação de Categoria Administrativa e Grande Região, os conceitos 2 e 4 foram atribuídos a cinco ( 25,0 \% ) e a três cursos ( 15,0 \% ), respectivamente.
Nenhum curso recebeu os conceitos 1 ou 5, e nenhum ficou Sem Conceito ( SC ).
As Instituições Públicas na região Sul participaram com 18 cursos, 47,4 \% dos cursos da região.
O conceito modal foi 4, atribuído a nove cursos ( 50,0 \% ).
Aos demais cursos foram atribuídos os conceitos 2 ( dois cursos, 11,1 \% ), 3 ( cinco cursos, 27,8 \% ) e 5 ( um curso, 5,6 \% ).
Nenhum curso da região Sul oferecido por IES Pública recebeu conceito 1.
Apenas um curso ( 5,6 \% ) ficou Sem Conceito ( SC ).
Na região Centro-Oeste, os dois únicos cursos participantes eram de Instituições Públicas ( 100,0 \% em termos regionais ).
Um desses cursos recebeu conceito 3, e o outro, conceito 4.
Sendo assim, nenhum curso ficou Sem Conceito ( SC ) ou recebeu os demais conceitos.
O Gráfico 5.2 apresenta a distribuição do Conceito Enade segundo a Categoria Administrativa da IES.
Os cursos em IES Públicas ( linha azul ) apresentam uma poligonal mais à direita do que os em IES Privadas ( linha verde ) e, consequentemente, uma distribuição de curso ( 25,0 \% ).
Nenhum curso ficou Sem Conceito ( SC ), assim como nenhum recebeu os demais conceitos nessa combinação de categoria e região.
Na região Nordeste, a rede privada concentrou seis dos 19 cursos participantes, o equivalente a 31,6 \% do total da região.
Dentre os cursos oferecidos por IES Privadas no Nordeste, cinco cursos ( 83,3 \% ) foram avaliados com conceito 2, o conceito modal.
Um curso obteve conceito 3 ( 16,7 \% ).
Nenhum curso ficou Sem Conceito ( SC ) ou com os demais conceitos nessa região.
As Instituições Públicas da região Nordeste participaram com 13 cursos ( 68,4 \% do total da região ).
Desses, sete obtiveram conceito 4 ( 53,8 \% ), o valor modal, seguidos por quatro cursos, que receberam conceito 3 ( 30,8 \% ), e por dois cursos, que receberam conceito 5 ( 15,4 \% ).
Na região Sudeste, a proporção de cursos da rede privada, 66,2 \%, foi mais elevada do que a registrada nas demais regiões brasileiras, correspondendo a 47 dos 71 cursos participantes.
Nessa categoria e região, o conceito modal foi 2, atribuído a 24 cursos ( 51,1 \% ).
Nenhum curso ficou Sem Conceito ( SC ) ou recebeu conceito 5.
Os demais foram avaliados com conceito 1 ( dois cursos, 4,3 \% ), conceito 3 ( 17 cursos, 36,2 \% ) e conceito 4 ( quatro cursos, 8,5 \% ).
Dentre os 24 cursos oferecidos por Instituições Públicas na região Sudeste, o conceito modal foi 4, atribuído a nove cursos ( 37,5 \% ).
Os demais cursos receberam os conceitos 2 ( um curso, 4,2 \% ), 3 ( sete cursos, 29,2 \% ) e 5 ( seis cursos, 25,0 \% ).
Na região Sudeste, nessa categoria, um curso ( 4,2 \% ) ficou Sem Conceito ( SC ), e nenhum curso recebeu conceito 1.
As Instituições Privadas concentraram 20 dos 38 cursos participantes da região Sul, 52,6 \% do total regional.
Desses, 12 ( 60,0 \% ) obtiveram conceito 3, o conceito modal.
Nessa combinação de Categoria Administrativa e Grande Região, os conceitos 2 e 4 foram atribuídos a cinco ( 25,0 \% ) e a três cursos ( 15,0 \% ), respectivamente.
Nenhum curso recebeu os conceitos 1 ou 5, e nenhum ficou Sem Conceito ( SC ).
As Instituições Públicas na região Sul participaram com 18 cursos, 47,4 \% dos cursos da região.
O conceito modal foi 4, atribuído a nove cursos ( 50,0 \% ).
Aos demais cursos foram atribuídos os conceitos 2 ( dois cursos, 11,1 \% ), 3 ( cinco cursos, 27,8 \% ) e 5 ( um curso, 5,6 \% ).
Nenhum curso da região Sul oferecido por IES Pública recebeu conceito 1.
Apenas um curso ( 5,6 \% ) ficou Sem Conceito ( SC ).
Na região Centro-Oeste, os dois únicos cursos participantes eram de Instituições Públicas ( 100,0 \% em termos regionais ).
Um desses cursos recebeu conceito 3, e o outro, conceito 4.
Sendo assim, nenhum curso ficou Sem Conceito ( SC ) ou recebeu os demais conceitos.
O Gráfico 5.2 apresenta a distribuição do Conceito Enade segundo a Categoria Administrativa da IES.
Os cursos em IES Públicas ( linha azul ) apresentam uma poligonal mais à direita do que os em IES Privadas ( linha verde ) e, consequentemente, uma distribuição de Conceitos Enade com valores maiores.
Corroborando essa informação, 100 \% dos cursos em IES Privadas receberam conceito 4 ou abaixo desse valor, já que a poligonal cumulativa alcança 100 \% no conceito 4.
Em contrapartida 84,7 \% dos cursos em IES Públicas receberam conceito 4 ou abaixo desse valor. 5.3 CONCEITOS POR ORGANIZAÇÃO ACADÊMICA E GRANDE REGIÃO Na Tabela 5.3, encontra -se a distribuição dos conceitos atribuídos aos cursos participantes do Enade/2017 na Área de Engenharia Química, por Organização Acadêmica, segundo as Grandes Regiões brasileiras.
Dos 135 cursos de Engenharia Química participantes, 91 eram oferecidos em Universidades ; 18, em Centros Universitários ; 24, em Faculdades ; e dois, em CEFET/IFET.
Essa distribuição corresponde a, respectivamente, 67,4 \%, 13,3 \%, 17,8 \% e 1,5 \% dos cursos.
De acordo com os dados apresentados, dos nove cursos avaliados com conceito 5, sete eram vinculados a Universidades.
Os cursos vinculados a esse tipo de Organização Acadêmica tiveram o conceito 3 como conceito modal, atribuído a 32 cursos ( 35,2 \% ).
Os demais cursos vinculados a Universidades receberam os conceitos 1 ( um curso, 1,1 \% ), 2 ( 18 cursos, 19,8 \% ), 4 ( 31 cursos, 34,1 \% ) e 5 ( sete cursos, 7,7 \%, como já mencionado ).
Ficaram Sem Conceito ( SC ) dois cursos ( 2,2 \% ).
Dentre os cursos em Centros Universitários, o conceito modal foi 2, atribuído a dez cursos ( 55,6 \% ).
Nenhum curso vinculado a esse tipo de Organização Acadêmica ficou Sem Conceito ( SC ) ou com conceitos 1 ou 5.
Os outros cursos receberam os conceitos 3 ( sete cursos, 38,9 \% ) e 4 ( um curso, 5,6 \% ).
Dos 24 cursos mantidos por Faculdades, dez ( 41,7 \% ) receberam o conceito 2 e outros dez ( 41,7 \% ) o conceito 3, os conceitos modais.
O segundo conceito mais frequentemente atribuído aos cursos oferecidos por essa categoria de IES foi o conceito 4, recebido por dois cursos ( 8,3 \% ).
Na sequência, os conceitos 1 e 5 foram atribuídos a um curso, cada um ( 4,2 \% ).
Nenhum curso ficou Sem Conceito ( SC ).
Dos dois cursos mantidos por CEFET/IFET, um ( 50,0 \% ) recebeu conceito 3, e outro ( 50,0 \% ) recebeu conceito 5.
Nenhum curso ficou Sem Conceito ( SC ) ou recebeu os demais conceitos.
Considerando -se separadamente as regiões brasileiras, verifica -se que, na região Norte, as Universidades participaram com quatro dos cinco cursos participantes ( 80,0 \% ).
Três cursos receberam conceito 3 ( 75,0 \% ), conceito modal, e um curso recebeu conceito 4 ( 25,0 \% ).
Sendo assim, nenhum curso ficou Sem Conceito ( SC ) ou recebeu os demais conceitos nessa combinação de categoria e região.
Os Centros Universitários da região Norte foram representados por apenas um curso, que recebeu conceito 2.
As Faculdades e os CEFET/IFET não participaram com cursos nessa região.
Na região Nordeste, as Universidades participaram com 14 dos 19 cursos da Área de Engenharia Química oferecidos na região ( 73,7 \% ), dos quais nenhum ficou Sem Conceito ( SC ) ou recebeu conceito 1.
O conceito modal foi 4, atribuído a sete cursos ( 50,0 \% ).
Os conceitos 2 e 5 foram atribuídos a um curso ( 7,1 \% ), cada um, e o conceito 3 foi atribuído a cinco cursos ( 35,7 \% ).
Os Centros Universitários contaram com dois cursos participantes na região Nordeste, que obtiveram o conceito 2.
Assim, nenhum curso oferecido por Centros Universitários no Nordeste ficou Sem Conceito ( SC ) ou recebeu os demais conceitos.
As Faculdades foram também representadas por dois cursos na região Nordeste, sendo que ambos, também, receberam conceito 2.
Já os CEFET/IFET participaram com apenas um curso, que recebeu conceito 5.
Na região Sudeste, as Universidades concentraram 39 dos 71 cursos de Engenharia Química ( 54,9 \% ).
Dentre os cursos em Universidades, nessa região, o conceito modal foi 2, atribuído a 12 cursos ( 30,8 \% ).
Os demais cursos receberam os conceitos 1 ( um curso, 2,6 \% ), 3 ( nove cursos, 23,1 \% ), 4 ( 11 cursos, 28,2 \% ) e 5 ( cinco cursos, 12,8 \% ).
Um curso ( 2,6 \% ) ficou Sem Conceito ( SC ).
Os Centros Universitários participaram com 12 cursos na região Sudeste ( 16,9 \% ), dos quais seis obtiveram o conceito modal 2 ( 50,0 \% ).
O conceito 3 foi atribuído a cinco cursos ( 41,7 \% ), e o conceito 4, a um curso ( 8,3 \% ).
Nenhum curso ficou Sem Conceito ( SC ) nessa combinação de organização acadêmica e região, assim como nenhum recebeu os conceitos 1 ou 5.
Já as Faculdades foram representadas por 19 cursos na região Sudeste ( 26,8 \% ), dos quais nenhum ficou Sem Conceito ( SC ).
Os conceitos 1, 4 e 5 foram atribuídos a um curso ( 5,3 \% ), cada um.
O conceito modal 3 foi atribuído a nove cursos ( 47,4 \% ), e o conceito 2, a sete cursos ( 36,8 \% ).
Os CEFET/IFET participaram com apenas um curso, que recebeu conceito 3.
Dos 38 cursos da região Sul, 32 ( 84,2 \% ) eram de Universidades, para os quais o conceito modal foi 3, atribuído a 14 cursos ( 43,8 \% ).
Nesse tipo de organização, um curso ( 3,1 \% ) ficou Sem Conceito ( SC ) e um curso recebeu conceito 5 ( 3,1 \% ).
Os demais cursos receberam os conceitos 2 e 3 ( cinco, 15,6 \%, e 11 cursos, 34,4 \%, respectivamente ).
Dos três cursos participantes de Centros Universitários da região Sul ( 7,9 \% ), nenhum ficou Sem Conceito ( SC ) ou recebeu os conceitos 1, 4 ou 5.
O conceito modal, 3, foi atribuído a dois cursos ( 66,7 \% ), e o outro curso recebeu conceito 2 ( 33,3 \% ).
Quanto aos três cursos vinculados a Faculdades na região Sul ( 7,9 \% ), os conceitos 2, 3 e 4 foram atribuídos, cada um, a um curso ( 33,3 \% ).
Nessa região não foram avaliados cursos em CEFET/IFET.
Na região Centro-Oeste, os dois únicos cursos participantes eram de Universidades, um dos cursos recebeu conceito 3, e o outro, conceito 4.
Centros Universitários, Faculdades e CEFET/IFET da região Centro-Oeste não tiveram cursos participantes.
O Gráfico 5.3 apresenta a distribuição cumulativa do Conceito Enade segundo a Organização Acadêmica da IES.
Os cursos em CEFET/IFET ( linha roxa ) apresentam uma poligonal mais à direita do que as demais, denotando uma melhor distribuição de Conceitos Enade.
No outro extremo, estão os cursos oferecidos em Centros Universitários ( linha verde ) com a poligonal mais à esquerda e, consequentemente, uma distribuição de Conceitos Enade com menores valores.
IES de Faculdades ( linha vermelha ) e de Universidades ( linha azul ) apresentam uma situação intermediária.

% matched lemmas: 
\end{textsample}
